% Independent written verification of the parent's full-radical group lead.
% Isolated candidate; not an assertion that K=0 in the geometric application.
\subsection{The exact normal closure of complete hyperplane radicals}

\begin{lemma}[Radicals generate a pointwise stabilizer]
\label{lem:full-radicals-pointwise-stabilizer}
\status{Candidate; elementary independent verification of the root's lead.}
Let \(V\) be a finite-dimensional vector space over \(\mathbf F_2\).
For \(0\ne\ell\in V^*\), put
\[
 U_\ell=\{1+v\otimes\ell:\ v\in\ker\ell\}.
\]
For a nonempty set \(\mathscr S\subset V^*\setminus\{0\}\), set
\(W=\langle\mathscr S\rangle\), \(K=W^\perp\), and
\(N=\langle U_\ell:\ell\in\mathscr S\rangle\).
Then
\[
             N=\{g\in\operatorname{GL}(V):g|_K=1\}.
                                                               \tag{NR1}
\]
The group on the right is the pointwise stabilizer, not only the
setwise stabilizer, of \(K\).
\end{lemma}
\begin{proof}
Every generator fixes \(K\), giving one inclusion.  Choose a basis
\(\ell_1,\ldots,\ell_a\) of \(W\) from \(\mathscr S\), and a
complement \(Q\) of \(K\), with basis \(q_1,\ldots,q_a\) dual to
those functionals.  For \(i\ne j\), the element
\(1+q_j\otimes\ell_i\) belongs to \(U_{\ell_i}\).  These elementary
matrices generate \(\operatorname{GL}(Q)\), because
\(\operatorname{GL}(Q)=\operatorname{SL}(Q)\) over \(\mathbf F_2\).
When \(a=1\), this factor is the trivial group, so the same assertion
still holds.  For every \(k\in K\), the element
\(1+k\otimes\ell_i\) also belongs to \(U_{\ell_i}\).  Such elements
generate all of \(\operatorname{Hom}(Q,K)\).  In the decomposition
\(V=K\oplus Q\), the pointwise stabilizer is exactly
\[
       \operatorname{Hom}(Q,K)\rtimes\operatorname{GL}(Q),
\]
so it is contained in \(N\), proving equality.
\end{proof}

\begin{proposition}[The full residual parabolic is determined]
\label{prop:normal-full-radical-parabolic}
\status{Candidate; does not exclude a nonzero residual space.}
Let \(G\le\operatorname{GL}(V)\) contain one complete radical
\(U_{\ell_0}\), and let \(N\) be its normal closure in \(G\).
Set
\[
 W=\langle\ell_0\circ g^{-1}:g\in G\rangle,\qquad K=W^\perp,
 \qquad H=\operatorname{im}(G\longrightarrow\operatorname{GL}(K)).
\]
Then \(K\) is \(G\)-stable, \(N\) is the full pointwise stabilizer
of \(K\) in \(\operatorname{GL}(V)\), and, after choosing a
complement of \(K\),
\[
       G=\operatorname{Hom}(V/K,K)\rtimes
          \bigl(\operatorname{GL}(V/K)\times H\bigr).             \tag{NR2}
\]
In particular \(K=0\) implies \(G=\operatorname{GL}(V)\).
If \((V^*)^G=0\), then \(\dim(V/K)\ge2\), but that condition
alone does not imply \(K=0\).
\end{proposition}
\begin{proof}
Conjugation gives
\(gU_{\ell_0}g^{-1}=U_{\ell_0\circ g^{-1}}\), so the preceding
lemma identifies \(N\).  The span \(W\), hence its annihilator
\(K\), is \(G\)-stable.  The restriction map from the setwise
stabilizer of \(K\) to \(\operatorname{GL}(K)\) has kernel
precisely \(N\).  Since \(G\) contains that whole kernel, it is
the entire inverse image of its image \(H\).  This proves (NR2)
and also justifies the displayed split form; no unproved
splitting of an arbitrary group extension is used.

If \(\dim W=1\), then its unique nonzero vector is fixed by \(G\),
because the ground field is \(\mathbf F_2\).  This contradicts
\((V^*)^G=0\).  If instead \(\dim W\ge2\), the pointwise
stabilizer \(N\) already has no fixed covector: its shear subgroup
forces such a covector to vanish on \(K\), and
\(\operatorname{GL}(V/K)\) has no nonzero fixed vector in
\((V/K)^*\).  Thus the absence of fixed covectors cannot rule
out a nonzero \(K\) in this case.
\end{proof}

\paragraph{Geometric scope.}
The all-odd Cartier theorem supplies the hypothesis that one
\(U_{\ell_0}\) is contained in the actual finite monodromy image.
The classification above isolates exactly what remains: the
annihilator \(K\) of the orbit of its normal may still be nonzero.
No boundary family, invariant-vector calculation, or global
connectedness assertion proving \(K=0\) is supplied here.
No numerical computation is required for either elementary proof.
